\documentclass[
a4paper, 
11pt
]{article}

%%%%%%%%%% Pacotes usados nesse documento %%%%%%%%%%
\usepackage[
lmargin = 3cm, 
tmargin = 3cm, 
rmargin = 2cm, 
bmargin = 2cm
]{geometry}
\usepackage{float} % Posicionamente de tabelas
\usepackage{amsmath} % Equações matemáticas
\usepackage{textcomp, gensymb} % Símbolos matemáticos, por exemplo \degree
\usepackage{multirow} % Para tabelas com células mescladas
\usepackage{array} % Tabelas com colunas de largura fixa 
\usepackage{enumitem} % Retirar identação do espaço \begin{itemize}
\usepackage{graphicx} % Utilizar imagens
\usepackage[brazil]{babel} % Texto automáticos serão gerados em português, por exemplo Tabela ao invés de Table.
\usepackage{datetime2} % Para obter o ano atual

% Formatação do texto no título de seções do documento
\usepackage[explicit]{titlesec}
\titleformat{\section}{\normalfont\Large\scshape}{}{0em}{#1\ \thesection}

% Definindo novas cores
\usepackage{color}
\definecolor{bblue}{rgb}{0.54,0.81,0.94}

% Utilizando ancoras de referência
\usepackage{hyperref}
\hypersetup{
colorlinks = true,
linkcolor = bblue,
anchorcolor = bblue,
citecolor = bblue,
filecolor = bblue,
urlcolor = bblue
}

 \DeclareMathOperator{\sen}{sen}

%%%%%%%%%% Início do documento %%%%%%%%%%
\begin{document}
\DTMsavenow{mytoday}

% Capa do Relatório
\begin{titlepage}
\begin{figure}
\centering

\includegraphics[width=3cm]{images/logo-ufc.png}
\end{figure}
\begin{center}
{\large \sc \noindent
{\LARGE Universidade Federal do Ceará} \\[0.5cm]
{Centro de Tecnologia} \\[0.5cm]
{Departamento de Engenharia de Teleinformática} \\[0.5cm]
{Introdução ao Processamento Digital de Imagens} \\[0.5cm]
{\DTMfetchday{mytoday}-\DTMfetchmonth{mytoday}-\DTMfetchyear{mytoday}} \\[4.0cm]
}
\end{center} 
\begin{center}
{\huge \sc \noindent Projeto 1 \\[0.5cm] Comparação das Filtragens nos Domínios Espacial e da Frequência} \\[3.0cm]
\end{center}
\begin{flushleft}
{\large \sc \noindent Aluno: Bruno Said Alves de Souza} \\[0.2cm]
{\large \sc \noindent Matrícula: 547860} \\[0.2cm]
{\large \sc \noindent Curso: Engenharia de Computação} \\[0.2cm]
{\large \sc \noindent Turma: TIO176 - T01} \\[0.2cm]
{\large \sc \noindent Professor: Paulo Cesar Cortez}
\end{flushleft}
\vspace*{\fill}
\begin{center}
{\large \sc \noindent
{Fortaleza} \\[0.2cm]
{\DTMfetchyear{mytoday}}
}
\end{center}
\end{titlepage}

\section{Discussão Técnica}
Podemos aplicar filtro as imagens tanto no domínio do espaço como no domínio da 
frequência. Além disso, podemos criar a mesma filtragem no domínio do espaço no
domínio da frequência usando a Tranformada de Fourier. Nesse projeto vamos criar
as menas images que serião criadas pelos filtro Laplaciano espacial. Para
realizar esse operação nos pricedemos da seguinte maneira.

\begin{itemize}
  \item Obter a imagem M x N que se deseja aplicar o filtro.
  \item Expandir a imagem para P x Q. Geralmente P = 2M e Q = 2N.
  \item Aplicar a Transformada de Fourier a imagem expandidad.
  \item Expandir o filtro no domínio do espaço para a mesma dimenção P x Q.
  \item Aplicar a Tranformada de Fourier ao filtro expandido.
  \item Multiplicar os dois filtro pela multiplicação de Hadamard, ou seja, 
  item por item.
  \item Aplicar a Transformada Inversa no resultado obtido
  \item Recordar a imagem M x N
  \item Resultado final encotrado
\end{itemize}

Todos esses paços foram aplicados na imagem a seguir.

\begin{figure}[ht]
  \centering
  \includegraphics[width = 15cm]{"./images/moons.png"}
\end{figure}

\section{Conclusões}
Podemos concluir que as imagens finais são identicas as geradas pelo filtro no
domínio espacial. Sendo assim, uma vantegem da troca de domínios é a velocidade
no precessamento da imagem pois qualquer um dos filtros pode ser obtido nos em
qualquer um dos dois domínios.

\end{document}
